\section*{Chapter \ref{ch:dv:managing-an-exotic-book} --- Managing an Exotic Book}

\begin{solution}{exo:dv:managing-an-exotic-book:1}
Its value depends materially on the correlation of the two indices over three years, an input no market quotes: an unobservable, Level 3 input. It could move to
Level 2 if the correlation became observable for substantially the note's remaining term (quoted \omterm{def:dv:multi-asset-options:corrswap}{correlation swaps} or baskets on the pair), or if the correlation's
effect on the value became insignificant.
\end{solution}

\begin{solution}{exo:dv:managing-an-exotic-book:2}
The firm owes the note. The higher the note's value, the lower the firm's book value, so the prudent point, the value it is 90\% sure it could exit at or better, is at
the high end of the plausible note values. With four models and the 90\% point below the first plotting position, it is the highest value: 98.62.
\end{solution}

\begin{solution}{exo:dv:managing-an-exotic-book:3}
$0.621+0.621+0.433+0.325=2.000$ million. Deferring everything net of bid--offer recognises nothing at inception and defers $2.000-0.621=1.379$ million.
\end{solution}

\begin{solution}{exo:dv:managing-an-exotic-book:4}
Not by construction: the \omterm{def:dv:managing-an-exotic-book:bidoffer}{bid--offer reserve} prices closing the \omterm{def:dv:greeks-and-the-hedging-pnl:greeks}{vega} that is known; the \omterm{def:dv:managing-an-exotic-book:model}{model reserve} prices not knowing which volatility is the right one. They coincide
because the alternative model here sits half a point of volatility away and the spread is half a point. A different model set, or a different spread, separates them.
The EU standard asks that model risk not include what the market-price reserve already covers, so the check belongs in the reserve's documentation.
\end{solution}

\begin{solution}{exo:dv:managing-an-exotic-book:5}
The note pays off on the worse of the two indices. With higher correlation the indices move together, the worse one is less far below the other, and knock-in and missed
coupons are less likely: the note is worth more to its holder.
\end{solution}

\begin{solution}{exo:dv:managing-an-exotic-book:6}
The firm is long \omterm{def:dv:greeks-and-the-hedging-pnl:greeks}{vega}: five points less volatility alone costs 5.89 million. A crash reaches the knock-in, where the firm's long put pays and outweighs the volatility loss
(7.02 million at $-30\%$). A moderate fall gains little at unchanged volatility (1.31 million at $-10\%$), and at lower volatility the knock-in is less likely still, so
nothing offsets the \omterm{def:dv:greeks-and-the-hedging-pnl:greeks}{vega} loss: $-6.31$ million at $-10\%$ with five points less volatility.
\end{solution}

\begin{solution}{exo:dv:managing-an-exotic-book:7}
Over 0.2--0.8 the note is worth 97.33 to 99.12; the 90\% point is 98.93 and the reserve 0.935. The deferred amount becomes $0.621+0.935=1.556$ and the recognised
\omterm{def:dv:managing-an-exotic-book:dayone}{day-one P\&L} $2.000-0.621-1.556=-0.177$ million: at that range, the sale at par does not cover its reserves.
\end{solution}

\begin{solution}{exo:dv:managing-an-exotic-book:8}
Adding every reserve at its 90\% point assumes that every position is exited at its worst at once. That is not conservative by design; it is miscalibrated, and it makes
the reserve useless as a measure of what the book could lose. The prudent-valuation standard's aggregation (half of the sum under its first method) exists for that
reason. Conversely, offsetting positions should be netted before reserving, or the calculation charges twice for one exposure.
\end{solution}

\begin{solution}{pb:dv:managing-an-exotic-book:1}
\textbf{1.} 7.95\% a year. The firm is long the investors' knock-in put (long \omterm{def:dv:greeks-and-the-hedging-pnl:greeks}{vega}, long skew) and short correlation.
\textbf{2.} Observable: the indices, rates, dividends, at-the-money volatilities. Not observable: the three-year correlation, and the choice of volatility for the knock-in.
\textbf{3.} Level 3. The day-one difference may be recognised only when fair value is evidenced by Level 1 prices or by observable data only; otherwise it is deferred.
\textbf{4.} \omterm{def:dv:greeks-and-the-hedging-pnl:greeks}{Vegas} $+0.533$ and $+0.693$ per point per 100; deltas $-0.143$ and $-0.279$ per 1\% before hedging (the firm is short the note).
\textbf{5.} The deltas, and part of the \omterm{def:dv:greeks-and-the-hedging-pnl:greeks}{vega} with listed options; not the correlation, not the long-dated skew.
\textbf{6.} 0.266 and 0.346 for the \omterm{def:dv:greeks-and-the-hedging-pnl:greeks}{vegas} and 0.008 for the deltas: 0.621 million (after rounding).
\textbf{7.} 98.00, 94.33, 98.62, 97.98; \omterm{def:dv:managing-an-exotic-book:model}{model reserve} 0.621 million.
\textbf{8.} 97.64 to 98.50; the 90\% point 98.43; reserve 0.433 million.
\textbf{9.} Half of $0.621+0.621+0.433$: 0.838 million.
\textbf{10.} 69 days at an assumed 100\,000 per point a day and a tenth of the volume: more than ten, so a concentration adjustment is due.
\textbf{11.} Charged 0.621 (bid--offer), deferred 1.054, recognised 0.325 million.
\textbf{12.} The correlation becomes observable within 0.47--0.53 and the note has two years left: the correlation reserve falls to 0.071 and the \omterm{def:dv:managing-an-exotic-book:model}{model reserve} to 0.559; 0.424
million is released.
\textbf{13.} 0.362 million.
\textbf{14.} With the indices at 80\% the note is worth 86.99 and nearer its barrier: the \omterm{def:dv:managing-an-exotic-book:model}{model reserve} rises to 0.649 and the correlation reserve is 0.202, so only 0.203
million is released.
\textbf{15.} It should protect a moderate fall with falling volatility: sell less of the upside hedge, or keep some short-dated \omterm{def:dv:greeks-and-the-hedging-pnl:greeks}{vega} as the indices fall, within the limits.
\textbf{16.} Slow in timing, conservative in amount: the profit comes back only when the uncertainty is resolved, and if the resolution goes badly, some of it never does.
\textbf{17.} It lowers the reward for margins that live in unobservable inputs, which is its purpose, and it pushes desks to trade structures whose inputs can be observed.
\textbf{18.} An independent valuation-control function, not the desk: the desk proposes, valuation control decides and documents.
\textbf{19.} Of a 2 million margin on 100 million of notes, 0.325 million is recognised at inception (0.621 charged, 1.054 deferred); 0.424 million is released after one year,
0.362 of it as the correlation becomes observable.
\textbf{20.} To value a position at what the firm could exit at, not at what its model says.
\end{solution}

\begin{solution}{iq:dv:managing-an-exotic-book:1}
The difference at inception between the transaction price and the fair value. Under IFRS 9 it is recognised at once only if the fair value is evidenced by a Level 1 price or
by a valuation using observable data only; otherwise it is deferred and released as factors market participants would price change, including time.
\iqlookfor{observability and deferral.}
\end{solution}

\begin{solution}{iq:dv:managing-an-exotic-book:2}
Price the exotic under a set of models that fit the same vanillas (local, stochastic, local-stochastic volatility; alternative calibrations), take the range, and reserve the
distance from the booked value to a prudent point (the 90\% point in the EU standard). Document the set, review it as markets change, and avoid double counting with the
market-price reserves.
\iqlookfor{model set, prudent point, independence.}
\end{solution}

\begin{solution}{iq:dv:managing-an-exotic-book:3}
Long volatility at long expiries, long skew, short correlation (worst-ofs), dividends, and large gamma near the barriers as they approach. They are hard to hedge because few
market participants take the other side and the exposures are shared by every issuer.
\iqlookfor{the \omterm{def:dv:greeks-and-the-hedging-pnl:greeks}{vega} and correlation positions and their one-way market.}
\end{solution}

\begin{solution}{iq:dv:managing-an-exotic-book:4}
A grid of index moves and volatility moves, with correlation and dividend shocks, full revaluation, and the hedges in place; plus the path scenarios where indices approach the
barriers. The feared cell here is a moderate fall with falling volatility, where the hedges lose and the long \omterm{def:dv:greeks-and-the-hedging-pnl:greeks}{vega} does too, not the crash.
\iqlookfor{full revaluation, joint shocks, the non-obvious worst case.}
\end{solution}

\begin{solution}{iq:dv:managing-an-exotic-book:5}
The market-data snapshot, the model and its version, the model set and plausible ranges with their justification, the positions, the code version, and the computed values at each
step, so that the same inputs reproduce the same reserve.
\iqlookfor{snapshot, versions, ranges and their rationale.}
\end{solution}

\begin{solution}{iq:dv:managing-an-exotic-book:6}
The parameter reserve on it shrinks to the market's bid--offer, and the deferred day-one profit tied to it is released; the mark moves to the observed value, which can be a gain
or a loss against the booked mid.
\iqlookfor{release of the deferral, and a mark-to-market effect.}
\end{solution}
